
Neural networks trained with different random seeds produce different test accuracies under otherwise identical conditions. This phenomenon is widely acknowledged, yet no validated measurement protocol exists for quantifying the magnitude of seed-based variance in simple neural networks. Research publications routinely report mean $\pm$ standard deviation across 3-5 random seeds without establishing whether this sample size is adequate or what variance magnitude should be expected.

This measurement gap has implications for uncertainty quantification research. Methods such as Monte Carlo Dropout, Bayesian neural networks, and deep ensembles aim to quantify predictive uncertainty, but without a baseline measurement of variance from seed initialization alone, it is unclear what portion of reported uncertainty reflects the method's contribution versus natural initialization stochasticity.

Prior work by Picard et al. (2021) examined random seed effects on CIFAR-10 ResNet-18 using 10,000 seeds, demonstrating that seed-dependent variance exists under deterministic training. However, their work did not establish baselines for simpler architectures or validate statistical protocols for practical sample sizes. Rajput and Kumar (2023) provided theoretical guidance recommending $N\geq 30$ for machine learning experiments, but did not empirically test this criterion for neural network variance estimation with bootstrap stability analysis.

This work addresses the following research questions: (1) What is the magnitude of test accuracy variance from random seed initialization in simple MLPs trained on MNIST and Fashion-MNIST? (2) Does the $N\geq 30$ sample size criterion provide stable variance estimates in this context? (3) Can the causal mechanism (seed $\rightarrow$ initial weights $\rightarrow$ training trajectories $\rightarrow$ final variance) be validated through hypothesis decomposition?

We measured variance using N=30 independent random seeds across four experimental conditions (2 architectures $\times$ 2 datasets) under full determinism. The experimental protocol enforced that random seed was the sole source of variance by controlling all other factors (optimizer, hyperparameters, batch order, dataset splits). Variance measurements were subjected to statistical significance tests and bootstrap stability analysis. The causal mechanism was validated through three hypotheses testing seed independence, trajectory divergence, and variance stability.

Results show task-dependent variance ranging from 0.04\% (MNIST) to 0.59\% (Fashion-MNIST, 2-layer). All conditions exhibited statistically significant non-zero variance (p < 0.05). Mechanism validation confirmed that different seeds produce independent initial weights (p < 0.000001), which converge to different local minima after training. Bootstrap analysis revealed that N=30 is sufficient for detecting non-zero variance but yields confidence interval widths of 93-110\%, exceeding the 50\% stability threshold proposed in prior work.

The remainder of this paper is organized as follows. Section 2 reviews related work in uncertainty quantification and reproducibility. Section 3 describes the experimental methodology and hypothesis decomposition approach. Section 4 presents the experimental design and procedures. Section 5 reports quantitative results for variance measurements and mechanism validation. Section 6 discusses limitations, implications, and future work. Section 7 concludes.
